% =====================================================================
% IEEE Conference Paper
% Topic: Model Predictive Control for Energy-Efficient Building HVAC
% =====================================================================
% !TeX spellcheck = en-US
% !TeX encoding = utf8
% !TeX program = pdflatex
% !BIB program = bibtex

\documentclass[conference,a4paper]{IEEEtran}[2015/08/26]

\usepackage{balance}
\usepackage{iftex}
\usepackage[hyphens]{url}
\makeatletter
\g@addto@macro{\UrlBreaks}{\UrlOrds}
\makeatother

\usepackage{graphicx}
\usepackage{amsmath}
\usepackage{amssymb}
\usepackage{bm}

\usepackage[dvipsnames,table]{xcolor}
\usepackage{booktabs}
\usepackage{multirow}
\usepackage{array}
\usepackage{paralist}

\usepackage{tikz}
\usetikzlibrary{shapes.geometric,arrows.meta,positioning,calc,fit,
                decorations.markings,patterns}
\usepackage{pgfplots}
\pgfplotsset{compat=1.17}
\usepgfplotslibrary{groupplots}

\usepackage{placeins}

\usepackage[autostyle=true]{csquotes}

\usepackage[%
  square,
  comma,
  numbers,
  sort&compress
]{natbib}
\renewcommand{\bibfont}{\normalfont\footnotesize}

\usepackage[group-minimum-digits=4,per-mode=fraction]{siunitx}

\usepackage{hyperref}
\hypersetup{
  hidelinks,
  colorlinks=true,
  allcolors=black,
  pdfstartview=Fit,
  breaklinks=true,
}

\usepackage[all]{hypcap}
\usepackage[caption=false,font=footnotesize]{subfig}
\usepackage[capitalise,nameinlink,noabbrev]{cleveref}

% ---- custom commands ----
\newcommand{\E}{\mathbb{E}}
\newcommand{\R}{\mathbb{R}}
\DeclareMathOperator*{\argmin}{arg\,min}
\newcommand{\vek}[1]{\bm{#1}}
\newcommand{\norm}[2]{\left\|#1\right\|_{#2}^{2}}

\hyphenation{op-tical net-works semi-conduc-tor pre-dic-tive}

% =====================================================================
\begin{document}
% =====================================================================

\title{Model Predictive Control for Energy-Efficient {HVAC} Management\\
       in Smart Building Automation Systems}

\author{%
  \IEEEauthorblockN{Zhenhan Zou}
  \IEEEauthorblockA{School of Automation\\
    Example University of Technology\\
    Shanghai, China\\
    \texttt{zouzhenhan@example.edu.cn}}
}

\maketitle

% =====================================================================
\begin{abstract}
% =====================================================================
Building heating, ventilation, and air-conditioning (HVAC) systems
account for approximately \SI{40}{\percent} of total building energy
consumption in commercial facilities. Classical thermostatic and
proportional-integral (PI) controllers maintain indoor temperature
setpoints without exploiting forecasts of weather, occupancy, or
electricity pricing, resulting in significant energy waste and avoidable
demand peaks. This paper presents a model predictive control (MPC)
framework for single-zone HVAC management that simultaneously minimizes
energy cost while satisfying thermal comfort constraints and equipment
capacity limits. The building thermal dynamics are described by a
second-order resistance-capacitance (RC) network model whose parameters
are identified from operational data. A constrained quadratic program
(QP) is solved at each sampling instant over a receding prediction
horizon of six hours, incorporating forecasts of outdoor temperature,
solar irradiance, occupancy-driven internal gains, and time-of-use (TOU)
electricity tariffs. A Luenberger disturbance observer augments the
state vector to compensate for unmodeled heat flows. The proposed
controller is evaluated on a high-fidelity simulation of a \SI{50}{\metre
\squared} office zone under both winter heating and summer cooling
scenarios. Compared with a rule-based thermostat and a PI controller,
MPC achieves \SI{31.2}{\percent} and \SI{17.5}{\percent} energy savings,
respectively, reduces thermal comfort violations by \SI{78.4}{\percent},
and cuts electricity cost by \SI{18.5}{\percent} by exploiting TOU
pricing for pre-conditioning strategies.
\end{abstract}

\IEEEpeerreviewmaketitle

% =====================================================================
\section{Introduction}
\label{sec:introduction}
% =====================================================================

The built environment is one of the largest consumers of primary energy
worldwide. In commercial buildings, HVAC systems typically account for
40--60\% of total electricity demand, making building climate control a
critical lever for reducing operational costs and carbon
emissions~\cite{ASHRAE2017,Drgona2020}. With the global penetration of
smart meters, building management systems (BMS), and Internet-of-Things
(IoT) sensors, it is now feasible to implement advanced control
algorithms that go far beyond the simple on/off logic of conventional
thermostats~\cite{Sturzenegger2016}.

Classical approaches to building climate control rely on rule-based
thermostats or proportional-integral-derivative (PID) controllers tuned
to maintain a fixed indoor temperature setpoint~\cite{Bengea2014}. While
these strategies are simple and reliable, they are inherently reactive:
they respond to current temperature deviations but cannot anticipate
future disturbances such as solar irradiance peaks, changes in occupancy,
or time-varying electricity prices. As a result, they frequently
over-heat or over-cool the building and miss opportunities to shift loads
toward cheaper off-peak hours~\cite{Halvgaard2012}.

Model predictive control (MPC) offers a systematic framework for
addressing these limitations. First proposed for industrial process
control by Richalet and Garcia in the late
1970s--1980s~\cite{Garcia1989}, MPC optimizes a finite-horizon cost
function subject to explicit process constraints at every time step,
then implements only the first element of the optimal control
sequence~\cite{Rawlings2017,Mayne2000}. In the building domain, the
prediction model captures thermal inertia, the cost function encodes
energy and comfort objectives, and the constraints enforce setpoint
comfort bands and equipment limits. This structure naturally supports
preview of weather forecasts, occupancy schedules, and electricity
tariffs~\cite{Oldewurtel2012}. Seminal demonstrations by Oldewurtel et
al.~\cite{Oldewurtel2012} and Sturzenegger et al.~\cite{Sturzenegger2016}
showed 17--41\% energy savings over rule-based baselines in Swiss office
buildings.

Despite its theoretical appeal, applying MPC to real buildings faces
several challenges. First, accurate thermal models are needed but are
hard to derive from first principles for complex geometries~\cite{Privara2013}.
Second, solving the optimization problem at each step requires embedded
computational resources. Third, disturbances such as unscheduled
occupancy or cloud cover introduce model mismatch that degrades
performance~\cite{Afram2014}. This paper addresses these challenges with
a computationally light RC-network model, an augmented disturbance
observer, and a warm-started active-set QP solver.

The main contributions of this paper are:
\begin{enumerate}
  \item A complete MPC design for single-zone HVAC control, from
    thermal modeling to constrained QP formulation and disturbance
    observer design, all in a unified state-space framework.
  \item A time-of-use electricity cost objective that enables
    pre-conditioning strategies, demonstrating economic benefits beyond
    pure energy minimization.
  \item A comprehensive simulation study comparing MPC against thermostat
    and PI baselines over a seven-day winter and a seven-day summer
    period, with detailed analysis of energy, comfort, and cost metrics.
  \item Sensitivity analysis showing robustness of the proposed
    controller to parametric uncertainty in the thermal model and
    forecast errors in weather data.
\end{enumerate}

The remainder of this paper is organized as follows.
\Cref{sec:related} reviews related work.
\Cref{sec:model} develops the building thermal model.
\Cref{sec:mpc} formulates the MPC optimization problem.
\Cref{sec:observer} presents the disturbance observer.
\Cref{sec:results} reports simulation results.
\Cref{sec:discussion} discusses findings, and
\cref{sec:conclusion} concludes the paper.

% =====================================================================
\section{Related Work}
\label{sec:related}
% =====================================================================

\subsection{Rule-Based and Classical Feedback Control}

The simplest building climate controllers are rule-based thermostats
that switch the HVAC system on and off when the temperature crosses
fixed hysteresis bounds. While inexpensive and robust, these produce
bang-bang power profiles that are poorly suited to demand-response
programs~\cite{Halvgaard2012}. PI and PID controllers~\cite{Bengea2014}
eliminate steady-state offset and provide smoother responses but are
still reactive: they carry no forecast information and cannot exploit
the thermal mass of the building structure as a free energy buffer.
Gain scheduling and fuzzy-logic variants~\cite{Afram2014} improve
part-load performance but require substantial manual tuning.

\subsection{Model Predictive Control for Buildings}

The application of MPC to HVAC control was pioneered by Ma et
al.~\cite{Ma2012}, who demonstrated significant energy savings on a
university campus cooling plant. Oldewurtel et al.~\cite{Oldewurtel2012}
conducted one of the largest empirical studies, covering 93 simulation
test cases for Swiss office buildings and reporting 17--41\% energy
savings. Sturzenegger et al.~\cite{Sturzenegger2016} extended this work
to a full-scale implementation at the ALSTOM building in Baden,
Switzerland, confirming 16\% savings compared with the existing BMS
controller. Morosan et al.~\cite{Morosan2010} developed a distributed
MPC architecture for multi-zone buildings, where each zone solves a
local subproblem and communicates boundary conditions with neighbors.
Privara et al.~\cite{Privara2013} emphasized the importance of model
identification as a prerequisite for building MPC and proposed a
subspace identification procedure. Bengea et al.~\cite{Bengea2014}
reported implementation experience in a mid-size commercial building,
highlighting the challenges of sensor integration and real-time solver
deployment. A comprehensive review of building MPC methods through 2020
is given by Drgo{\v{n}}a et al.~\cite{Drgona2020}.

\subsection{Economic and Demand-Response MPC}

Halvgaard et al.~\cite{Halvgaard2012} incorporated time-of-use
electricity prices directly into the MPC cost function, showing that
smart pre-cooling can shift \SI{30}{\percent} of cooling energy to
cheaper nighttime slots. Oldewurtel et al.~\cite{Oldewurtel2010}
proposed a stochastic MPC formulation that handles weather forecast
uncertainty via scenario trees. Parisio et al.~\cite{Parisio2014}
applied economic MPC to microgrid energy management, a closely related
problem. Fiorentini et al.~\cite{Fiorentini2019} combined MPC with
on-site photovoltaic generation and thermal storage, demonstrating
bill reductions of up to \SI{25}{\percent}. V{\'a}zquez-Canteli and
Nagy~\cite{Vazquez2019} reviewed reinforcement learning approaches
to demand response as an alternative to MPC when process models are
unavailable. Yang et al.~\cite{Yang2015} combined MPC with a
machine-learning prediction model, achieving robust performance under
changing occupancy patterns.

\subsection{Thermal Modeling for Control}

Reliable building thermal models are essential for MPC performance.
RC-network (lumped-parameter) models provide a good balance between
accuracy and tractability for control design~\cite{Bacher2011}.
Bacher and Madsen~\cite{Bacher2011} developed a systematic
identification methodology for continuous-time stochastic RC models
using data from occupied buildings. More detailed models based on
finite-element heat transfer or EnergyPlus are accurate but too
expensive to embed in real-time optimization~\cite{Underwood2014}.
Aswani et al.~\cite{Aswani2012} proposed a data-driven learning
approach that guarantees provably safe and robust MPC performance even
when the model is imperfect, providing a theoretical bridge between
identification accuracy and closed-loop stability.

% =====================================================================
\section{Building Thermal Model}
\label{sec:model}
% =====================================================================

\subsection{Second-Order RC Network}

We model a single thermal zone using a second-order resistance-capacitance
(RC) network, as illustrated in the left portion of \cref{fig:system}.
The network contains two thermal capacitances: the zone air mass
$C_a$ and the lumped building structure (walls, ceiling, floor) $C_s$.
The continuous-time dynamics are:
\begin{align}
  C_a \dot{T}_a &= \frac{T_s - T_a}{R_{as}}
                 + \frac{T_{\text{out}} - T_a}{R_{ao}}
                 + \dot{Q}_{\text{hvac}}
                 + \dot{Q}_{\text{int}},
  \label{eq:Ta}\\[4pt]
  C_s \dot{T}_s &= \frac{T_a - T_s}{R_{as}}
                 + \frac{T_{\text{out}} - T_s}{R_{so}}
                 + \alpha_s\, A_w\, I_{\text{sol}},
  \label{eq:Ts}
\end{align}
where $T_a$ is the zone air temperature (\si{\celsius}), $T_s$ is the
structural temperature (\si{\celsius}), $T_{\text{out}}$ is the outdoor
air temperature (\si{\celsius}), $R_{as}$ is the air-structure thermal
resistance (\si{\kelvin\per\watt}), $R_{ao}$ and $R_{so}$ are the
air-outdoor and structure-outdoor thermal resistances, $\dot{Q}_{\text{hvac}}$
is the HVAC thermal power output (\si{\watt}), $\dot{Q}_{\text{int}}$
represents internal heat gains from occupants and equipment (\si{\watt}),
$\alpha_s$ is the solar absorptivity of the structure, $A_w$ is the
effective window area (\si{\metre\squared}), and $I_{\text{sol}}$ is
solar irradiance (\si{\watt\per\metre\squared}).

\subsection{State-Space Formulation}

Defining the state vector $\vek{x} = [T_a,\; T_s]^\top \in \R^2$,
the control input $u = \dot{Q}_{\text{hvac}} \in \R$ (positive for
heating, negative for cooling), and the disturbance vector
$\vek{d} = [T_{\text{out}},\; \dot{Q}_{\text{int}},\; I_{\text{sol}}]^\top \in \R^3$,
\cref{eq:Ta,eq:Ts} may be written in continuous-time state-space form:
\begin{equation}
  \dot{\vek{x}} = A_c\, \vek{x} + B_c\, u + E_c\, \vek{d},
  \label{eq:ss_cont}
\end{equation}
where
\begin{align}
  A_c &= \begin{bmatrix}
           -\!\left(\tfrac{1}{C_a R_{as}} + \tfrac{1}{C_a R_{ao}}\right) & \tfrac{1}{C_a R_{as}} \\[4pt]
           \tfrac{1}{C_s R_{as}} & -\!\left(\tfrac{1}{C_s R_{as}} + \tfrac{1}{C_s R_{so}}\right)
         \end{bmatrix},
  \label{eq:Ac}
\end{align}
$B_c = [1/C_a,\; 0]^\top$, and $E_c$ contains the disturbance coupling
coefficients. Zero-order-hold discretization with sampling period
$T_s = \SI{900}{\second}$ (\SI{15}{min}) yields:
\begin{equation}
  \vek{x}_{k+1} = A\,\vek{x}_k + B\,u_k + E\,\vek{d}_k + \vek{w}_k,
  \label{eq:ss_disc}
\end{equation}
where $\vek{w}_k \sim \mathcal{N}(\vek{0}, Q_w)$ captures unmodeled
dynamics and $A = e^{A_c T_s}$, $B = A_c^{-1}(A - I)B_c$,
$E = A_c^{-1}(A - I)E_c$.

The output (measured zone air temperature) is:
\begin{equation}
  y_k = C\,\vek{x}_k + v_k, \quad
  C = \begin{bmatrix} 1 & 0 \end{bmatrix},
  \label{eq:output}
\end{equation}
where $v_k \sim \mathcal{N}(0, \sigma_v^2)$ is sensor noise.

\subsection{Model Identification}

The six RC parameters $\{C_a, C_s, R_{as}, R_{ao}, R_{so}, \alpha_s\}$
are identified by minimizing the one-step-ahead prediction error on
historical data using the Nelder-Mead simplex algorithm. Physical bounds
are imposed on each parameter based on material properties and building
geometry, regularizing the identification problem~\cite{Bacher2011}.
\Cref{tab:params} lists the identified parameter values used in all
simulations.

% =====================================================================
\section{MPC Controller Design}
\label{sec:mpc}
% =====================================================================

\subsection{Prediction Model and Horizon}

At time step $k$, the MPC controller solves a QP over a prediction
horizon of $N = 24$ steps (corresponding to \SI{6}{\hour} ahead).
Using~\cref{eq:ss_disc}, the predicted state trajectory is:
\begin{equation}
  \vek{x}_{k+i+1} = A\,\vek{x}_{k+i} + B\,u_{k+i} + E\,\hat{\vek{d}}_{k+i},
  \quad i = 0, \ldots, N-1,
  \label{eq:pred}
\end{equation}
where $\hat{\vek{d}}_{k+i}$ is the disturbance forecast at step $k+i$,
obtained from a weather service (outdoor temperature and solar irradiance)
and an occupancy schedule (internal gains).

\subsection{Objective Function}

The MPC cost function balances three objectives: setpoint tracking, TOU
energy cost, and penalization of large input changes:
\begin{equation}
  J = \sum_{i=0}^{N-1} \Bigl[
        q\bigl(T_{a,k+i} - T^{\text{ref}}\bigr)^2
        + c_{k+i}\, P_{k+i}
        + \lambda\, (\Delta u_{k+i})^2
      \Bigr],
  \label{eq:cost}
\end{equation}
where $q > 0$ is the thermal comfort weight, $T^{\text{ref}}$ is the
occupied setpoint, $c_{k+i}$ is the electricity tariff at step $k+i$
(\si{\euro\per\kilo\watt\hour}), $P_{k+i}$ is the electrical power
drawn by the HVAC unit, and $\lambda > 0$ penalizes abrupt control
changes. The electrical power consumed by the HVAC unit is related to
the thermal power output via the coefficient of performance:
\begin{equation}
  P_k = \frac{|u_k|}{\mathrm{COP}(u_k)},
  \label{eq:cop}
\end{equation}
where $\mathrm{COP}(u_k) = \mathrm{COP}_h$ for heating ($u_k > 0$) and
$\mathrm{COP}(u_k) = \mathrm{COP}_c$ for cooling ($u_k < 0$).

\subsection{Constraints}

The optimization is subject to the following constraints.

\textbf{Comfort constraints.} During occupied hours the zone temperature
must remain within the ASHRAE\,55 comfort band~\cite{ASHRAE55}:
\begin{equation}
  T^{\min}_{k+i} \le T_{a,k+i} \le T^{\max}_{k+i},
  \quad i = 1, \ldots, N,
  \label{eq:comfort}
\end{equation}
where $T^{\min} = \SI{20}{\celsius}$, $T^{\max} = \SI{26}{\celsius}$
during occupied periods (\SI{08:00}--\SI{18:00}) and
$T^{\min} = \SI{16}{\celsius}$, $T^{\max} = \SI{28}{\celsius}$ otherwise.

\textbf{Equipment capacity constraints.}
\begin{equation}
  u^{\min} \le u_{k+i} \le u^{\max},
  \quad i = 0, \ldots, N-1,
  \label{eq:cap}
\end{equation}
where $u^{\max} = \SI{6}{\kilo\watt}$ (full heating) and
$u^{\min} = -\SI{5}{\kilo\watt}$ (full cooling).

\textbf{Slew-rate constraint.}
\begin{equation}
  |\Delta u_{k+i}| \le \Delta u^{\max} = \SI{1}{\kilo\watt},
  \quad i = 0, \ldots, N-1.
  \label{eq:slew}
\end{equation}

\subsection{QP Formulation}

Stacking the predicted states and inputs into vectors
$\vek{X} = [\vek{x}_{k+1}^\top, \ldots, \vek{x}_{k+N}^\top]^\top$ and
$\vek{U} = [u_k, \ldots, u_{k+N-1}]^\top$, and substituting the
prediction model~\eqref{eq:pred} recursively yields the compact QP:
\begin{equation}
  \min_{\vek{U}}\quad \frac{1}{2}\vek{U}^\top H\,\vek{U} + \vek{f}^\top \vek{U}
  \quad \text{s.t.} \quad G\,\vek{U} \le \vek{h},
  \label{eq:qp}
\end{equation}
where $H \succ 0$ and $\vek{f}$ incorporate the prediction matrices and
cost weights, and the inequality $G\,\vek{U} \le \vek{h}$ encodes
constraints~\eqref{eq:comfort}--\eqref{eq:slew} after substituting the
disturbance forecasts. The positive definiteness of $H$ guarantees a
unique global optimum. Only $u_k$ (the first element of $\vek{U}^*$)
is applied to the plant; the horizon shifts by one step and the
optimization is repeated at $k+1$, implementing the receding-horizon
principle~\cite{Rawlings2017}.

% =====================================================================
\section{Disturbance Observer}
\label{sec:observer}
% =====================================================================

\subsection{Augmented State Model}

Real buildings are subject to unmodeled heat sources (e.g., unexpected
equipment loads, infiltration, thermal bridges) that bias the prediction
and degrade MPC performance. We model these as an additive disturbance
$d_k^a$ entering through the air-node equation:
\begin{equation}
  \vek{x}_{k+1}^{\text{aug}} =
    \underbrace{\begin{bmatrix} A & B_d \\ 0 & I \end{bmatrix}}_{A^{\text{aug}}}
    \vek{x}_k^{\text{aug}} +
    \underbrace{\begin{bmatrix} B \\ 0 \end{bmatrix}}_{B^{\text{aug}}}
    u_k +
    \underbrace{\begin{bmatrix} E \\ 0 \end{bmatrix}}_{E^{\text{aug}}}
    \vek{d}_k,
  \label{eq:aug}
\end{equation}
where $\vek{x}_k^{\text{aug}} = [\vek{x}_k^\top,\; d_k^a]^\top$ and
$B_d = [1/C_a,\; 0]^\top$. The augmented output equation is:
\begin{equation}
  y_k = C^{\text{aug}}\,\vek{x}_k^{\text{aug}} + v_k,
  \quad C^{\text{aug}} = [C \;\; 0].
  \label{eq:aug_out}
\end{equation}

\subsection{Luenberger Observer}

A Luenberger observer estimates the augmented state:
\begin{align}
  \hat{\vek{x}}_{k+1}^{\text{aug}} &=
    A^{\text{aug}}\hat{\vek{x}}_k^{\text{aug}}
    + B^{\text{aug}} u_k
    + E^{\text{aug}} \hat{\vek{d}}_k
    + L\bigl(y_k - C^{\text{aug}}\hat{\vek{x}}_k^{\text{aug}}\bigr),
  \label{eq:observer}
\end{align}
where the observer gain $L \in \R^{3 \times 1}$ is designed by pole
placement to achieve fast disturbance estimation with a time constant
of approximately $2\,T_s$. The estimated additive disturbance
$\hat{d}_k^a$ is then included as a constant over the prediction
horizon, improving prediction accuracy under slowly varying model
mismatch~\cite{Anderson1979,Simon2010}.

% =====================================================================
\section{Simulation Results}
\label{sec:results}
% =====================================================================

\subsection{Simulation Setup}

All simulations are conducted in MATLAB/Simulink using a higher-fidelity
two-zone building model (different from the RC model used for MPC
prediction) to represent the true plant, introducing a realistic
model-mismatch scenario. The QP~\eqref{eq:qp} is solved using
YALMIP~\cite{Lofberg2004} with QUADPROG as the back-end solver.
\Cref{tab:params} summarizes the identified thermal parameters; the MPC
design parameters are listed in \cref{tab:mpc_params}. Two evaluation
scenarios are considered: a seven-day winter period (average outdoor
temperature $-2$ to $\SI{8}{\celsius}$) and a seven-day summer period
(average outdoor $28$ to $\SI{38}{\celsius}$). Three controllers are
compared:
\begin{compactitem}
  \item \textbf{Thermostat:} hysteresis bang-bang, $\pm\SI{1}{\kelvin}$ band;
  \item \textbf{PI:} integral gain tuned for a \SI{5}{\minute} rise time;
  \item \textbf{MPC (proposed):} formulation of \cref{sec:mpc,sec:observer}.
\end{compactitem}

\begin{table*}[t]
  \caption{Identified Building Zone Thermal Parameters}
  \label{tab:params}
  \centering
  \begin{tabular}{lllS[table-format=2.3]l}
    \toprule
    \textbf{Parameter} & \textbf{Symbol} & \textbf{Description}
      & {\textbf{Value}} & \textbf{Unit} \\
    \midrule
    Air thermal mass         & $C_a$       & Lumped air + furniture capacity
      & 0.182  & \si{\mega\joule\per\kelvin} \\
    Structure thermal mass   & $C_s$       & Walls, ceiling, floor capacity
      & 2.640  & \si{\mega\joule\per\kelvin} \\
    Air-structure resistance & $R_{as}$    & Convective coupling
      & 0.052  & \si{\kelvin\per\watt} \\
    Air-outdoor resistance   & $R_{ao}$    & Infiltration + ventilation
      & 0.147  & \si{\kelvin\per\watt} \\
    Structure-outdoor resist.& $R_{so}$    & Wall conduction
      & 0.214  & \si{\kelvin\per\watt} \\
    Solar absorptivity       & $\alpha_s$  & Window solar transmission
      & 0.610  & --- \\
    Effective window area    & $A_w$       & South-facing glazing
      & 8.000  & \si{\metre\squared} \\
    Heating COP              & $\mathrm{COP}_h$ & Heat pump heating mode
      & 4.000  & --- \\
    Cooling COP              & $\mathrm{COP}_c$ & Heat pump cooling mode
      & 3.500  & --- \\
    HVAC rated capacity      & ---         & Max heating / cooling power
      & 6.000  & \si{\kilo\watt} \\
    \bottomrule
  \end{tabular}
\end{table*}

\begin{table}[t]
  \caption{MPC Design Parameters}
  \label{tab:mpc_params}
  \centering
  \begin{tabular}{llS[table-format=3.3]l}
    \toprule
    \textbf{Parameter} & \textbf{Symbol}
      & {\textbf{Value}} & \textbf{Unit} \\
    \midrule
    Sampling period       & $T_s$       & 900    & \si{\second} \\
    Prediction horizon    & $N$         & 24     & steps \\
    Comfort weight        & $q$         & 100    & \si{\watt\per\kelvin\squared} \\
    Input change weight   & $\lambda$   & 0.10   & \si{\watt\per\watt\squared} \\
    Occupied temp.\ min   & $T^{\min}$  & 20     & \si{\celsius} \\
    Occupied temp.\ max   & $T^{\max}$  & 26     & \si{\celsius} \\
    Max heating power     & $u^{\max}$  & 6.0    & \si{\kilo\watt} \\
    Max cooling power     & $|u^{\min}|$& 5.0    & \si{\kilo\watt} \\
    Max slew rate         & $\Delta u^{\max}$ & 1.0 & \si{\kilo\watt} \\
    Observer poles        & ---         & 0.200  & (discrete) \\
    \bottomrule
  \end{tabular}
\end{table}

\subsection{Temperature Tracking Performance}

\Cref{fig:temp} shows the zone air temperature over a representative
winter day under all three controllers. The thermostat produces the
characteristic bang-bang oscillation, remaining within the comfort band
but with large temperature swings. The PI controller overshoots upon
the morning step change from the night setback (\SI{18}{\celsius}) to
the occupied setpoint (\SI{22}{\celsius}) before settling. The MPC
controller pre-heats the zone starting at \SI{06:30} by exploiting the
thermal mass forecast, reaching the occupied setpoint smoothly by
\SI{08:00} with no overshoot. The pre-heating energy is drawn during the
off-peak tariff period (before 07:00), reducing cost. At \SI{18:00},
MPC allows the temperature to drift toward the night setback limit
rather than actively cooling, further saving energy.

\begin{figure*}[t]
  \centering
  \begin{tikzpicture}
    \begin{groupplot}[
      group style={
        group size=1 by 2,
        vertical sep=1.4cm,
      },
      width=\textwidth, height=4.5cm,
      grid=major, grid style={dotted,gray!50},
      tick label style={font=\scriptsize},
      label style={font=\footnotesize},
      every axis plot/.append style={thick},
      xmin=0, xmax=24,
      xtick={0,2,4,6,8,10,12,14,16,18,20,22,24},
      xticklabels={00:00,02:00,04:00,06:00,08:00,10:00,12:00,
                   14:00,16:00,18:00,20:00,22:00,24:00},
      x tick label style={font=\scriptsize, rotate=30, anchor=east},
    ]
    % ---------- top subplot: indoor temperatures ----------
    \nextgroupplot[
      ylabel={Temperature (\si{\celsius})},
      ymin=15, ymax=25,
      legend pos=north west,
      legend style={font=\scriptsize,fill=white,fill opacity=0.9,
                    draw opacity=1,text opacity=1},
    ]
    % Comfort band (occupied 8-18)
    \addplot[fill=green!15, draw=none, forget plot]
      coordinates{(8,20)(18,20)(18,26)(8,26)} -- cycle;

    % Setpoint step
    \addplot[black, dashed, thin, forget plot]
      coordinates{(0,18)(8,18)(8,22)(18,22)(18,18)(24,18)};

    % Thermostat
    \addplot[color=red!70!black] coordinates{
      (0,18.2)(1,17.8)(2,18.3)(3,17.7)(4,18.2)(5,17.9)
      (6,18.1)(7,17.8)(8,18.0)(9,22.5)(10,21.5)(11,22.5)
      (12,21.4)(13,22.5)(14,21.5)(15,22.4)(16,21.6)(17,22.5)
      (18,21.8)(19,18.4)(20,17.9)(21,18.3)(22,17.8)(23,18.2)(24,18.0)
    };
    \addlegendentry{Thermostat}

    % PI
    \addplot[color=blue!70!black] coordinates{
      (0,18.0)(1,18.0)(2,18.0)(3,18.0)(4,18.0)(5,18.0)
      (6,18.0)(7,18.0)(8,18.0)(9,23.2)(10,22.4)(11,22.1)
      (12,22.0)(13,22.1)(14,22.0)(15,22.1)(16,22.0)(17,22.1)
      (18,22.0)(19,17.2)(20,18.3)(21,18.1)(22,18.0)(23,18.0)(24,18.0)
    };
    \addlegendentry{PI controller}

    % MPC
    \addplot[color=orange!80!black] coordinates{
      (0,18.0)(1,18.0)(2,18.0)(3,18.0)(4,18.0)(5,18.0)
      (6,18.3)(7,20.6)(8,22.0)(9,22.1)(10,22.0)(11,22.1)
      (12,22.0)(13,22.1)(14,22.0)(15,22.0)(16,22.0)(17,21.9)
      (18,21.5)(19,20.2)(20,19.1)(21,18.5)(22,18.2)(23,18.0)(24,18.0)
    };
    \addlegendentry{MPC (proposed)}

    % ---------- bottom subplot: electrical power ----------
    \nextgroupplot[
      ylabel={HVAC Power (\si{\kilo\watt})},
      xlabel={Time of Day},
      ymin=-0.5, ymax=5.5,
    ]

    % Thermostat power
    \addplot[color=red!70!black, const plot] coordinates{
      (0,1.3)(1,0)(2,1.3)(3,0)(4,1.3)(5,0)(6,1.3)(7,0)
      (8,3.5)(9,0)(10,3.5)(11,0)(12,3.5)(13,0)(14,3.5)(15,0)
      (16,3.5)(17,0)(18,3.5)(19,0)(20,1.3)(21,0)(22,1.3)(23,0)(24,0)
    };

    % PI power
    \addplot[color=blue!70!black] coordinates{
      (0,0)(1,0)(2,0)(3,0)(4,0)(5,0)
      (6,0)(7,0)(8,5.0)(9,3.2)(10,2.1)(11,1.5)
      (12,1.5)(13,1.5)(14,1.5)(15,1.5)(16,1.5)(17,1.5)
      (18,1.5)(19,0)(20,0.8)(21,0.2)(22,0)(23,0)(24,0)
    };

    % MPC power
    \addplot[color=orange!80!black] coordinates{
      (0,0.5)(1,0.4)(2,0.3)(3,0.2)(4,0.2)(5,0.3)
      (6,0.8)(7,2.5)(8,1.2)(9,0.9)(10,0.7)(11,0.6)
      (12,0.5)(13,0.6)(14,0.5)(15,0.5)(16,0.4)(17,0.3)
      (18,0.2)(19,0.1)(20,0)(21,0)(22,0)(23,0)(24,0)
    };
    \end{groupplot}
  \end{tikzpicture}
  \caption{Top: zone air temperature over a representative winter day.
    Green shading marks the occupied comfort band (\SI{20}--\SI{26}{\celsius},
    08:00--18:00); dashed line shows the setpoint schedule.
    Bottom: corresponding HVAC electrical power consumption.
    MPC pre-heats during off-peak hours and avoids the large morning
    power spike of PI and thermostat control.}
  \label{fig:temp}
\end{figure*}

\subsection{Seven-Day Performance Comparison}

\Cref{tab:perf} summarizes the aggregate performance metrics over
seven-day winter and summer simulation periods. Thermal comfort
violations are counted as the number of \SI{15}{\minute} intervals in
which the zone temperature falls outside the occupied comfort band.
Energy is the total electrical energy consumed by the HVAC unit
(\si{\kilo\watt\hour}). Electricity cost uses a two-period TOU tariff:
\SI{0.12}{\euro\per\kilo\watt\hour} (off-peak, 22:00--07:00) and
\SI{0.28}{\euro\per\kilo\watt\hour} (peak, 07:00--22:00).

MPC reduces winter energy consumption from \SI{152.3}{\kilo\watt\hour}
(thermostat) to \SI{104.7}{\kilo\watt\hour}---a \SI{31.2}{\percent}
saving---and from \SI{143.7}{\kilo\watt\hour} (PI) by \SI{17.5}{\percent}.
Summer savings are \SI{29.8}{\percent} and \SI{16.1}{\percent} versus
thermostat and PI, respectively. Comfort violations are reduced by
\SI{78.4}{\percent} versus PI (from 37 to 8 intervals over both seasons).
Total TOU cost savings versus thermostat exceed \SI{18}{\percent}.

\begin{table*}[t]
  \caption{Seven-Day Performance Comparison (Winter and Summer)}
  \label{tab:perf}
  \centering
  \setlength{\tabcolsep}{5pt}
  \begin{tabular}{l
                  S[table-format=3.1] S[table-format=2.0] S[table-format=2.2]
                  S[table-format=3.1] S[table-format=2.0] S[table-format=2.2]}
    \toprule
    & \multicolumn{3}{c}{\textbf{Winter (7 days)}}
    & \multicolumn{3}{c}{\textbf{Summer (7 days)}} \\
    \cmidrule(lr){2-4}\cmidrule(lr){5-7}
    \textbf{Controller}
      & {Energy (kWh)} & {Comfort viol.} & {Cost (\euro)}
      & {Energy (kWh)} & {Comfort viol.} & {Cost (\euro)} \\
    \midrule
    Thermostat
      & 152.3 & 21 & 30.64
      & 187.4 & 26 & 38.22 \\
    PI controller
      & 143.7 & 19 & 27.81
      & 178.2 & 18 & 34.57 \\
    \textbf{MPC (proposed)}
      & \textbf{104.7} & \textbf{4} & \textbf{18.38}
      & \textbf{131.6} & \textbf{4} & \textbf{22.71} \\
    \midrule
    MPC vs.\ thermostat (\%)
      & {$-$31.2} & {$-$81.0} & {$-$40.0}
      & {$-$29.8} & {$-$84.6} & {$-$40.6} \\
    MPC vs.\ PI (\%)
      & {$-$17.5} & {$-$78.9} & {$-$33.9}
      & {$-$16.1} & {$-$77.8} & {$-$34.3} \\
    \bottomrule
  \end{tabular}
\end{table*}

\subsection{System Architecture}

\Cref{fig:system} shows the overall closed-loop architecture.
The MPC block receives the disturbance observer estimate, the weather
forecast, and the occupancy/tariff schedule, solves the QP, and
sends the optimal thermal power command to the HVAC unit. The observer
runs in parallel, updating the estimated additive disturbance at every
sampling instant.

\begin{figure}[t]
  \centering
  \begin{tikzpicture}[
    block/.style={rectangle, draw, fill=blue!8, rounded corners=2pt,
                  text width=2.0cm, align=center, minimum height=0.8cm,
                  font=\footnotesize},
    obsblock/.style={rectangle, draw, fill=orange!12, rounded corners=2pt,
                     text width=2.0cm, align=center, minimum height=0.8cm,
                     font=\footnotesize},
    foreblock/.style={rectangle, draw, fill=green!10, rounded corners=2pt,
                      text width=2.0cm, align=center, minimum height=0.8cm,
                      font=\footnotesize},
    arr/.style={-{Latex[length=2.5pt]}, thick},
    node distance=0.65cm and 0.85cm,
  ]
    % Main loop nodes
    \node[block]  (mpc)  {MPC\\Optimizer};
    \node[block, right=1.1cm of mpc]  (hvac) {HVAC\\Unit};
    \node[block, right=1.1cm of hvac] (bldg) {Building\\Zone};

    % Observer and forecast
    \node[obsblock, below=0.9cm of mpc]  (obs)  {Disturbance\\Observer};
    \node[foreblock, above=0.9cm of mpc] (fore) {Forecast\\$\hat{\vek{d}}_{k{:}k{+}N}$};

    % Arrows main path
    \draw[arr] (mpc)  -- node[above,font=\scriptsize]{$u_k^*$} (hvac);
    \draw[arr] (hvac) -- node[above,font=\scriptsize]{$\dot{Q}$} (bldg);

    % Feedback
    \coordinate (out) at ($(bldg.east)+(0.4,0)$);
    \draw[arr] (out) node[right,font=\scriptsize]{$y_k$}
               -- ++(0,-1.6) -| (obs.east);
    \draw[arr] (out) -- ++(0,1.6) -| (fore.east);

    % Observer to MPC
    \draw[arr] (obs.west) -- ++(-0.5,0) |-
               node[near end, left, font=\scriptsize]{$\hat{d}^a$}
               (mpc.south);

    % Forecast to MPC
    \draw[arr] (fore.west) -- ++(-0.5,0) |-
               node[near end, left, font=\scriptsize]{}
               (mpc.north);

    % Setpoint / tariff input
    \node[left=0.7cm of mpc, align=right, font=\scriptsize] (sp)
      {$T^{\text{ref}}$, $c_k$};
    \draw[arr] (sp) -- (mpc);

  \end{tikzpicture}
  \caption{Closed-loop MPC building control architecture. The MPC
    optimizer receives forecasts and observer estimates, solves the
    constrained QP, and sends the optimal thermal power command to the
    HVAC unit. The disturbance observer runs in parallel to compensate
    for unmodeled heat flows.}
  \label{fig:system}
\end{figure}

\subsection{Robustness to Forecast Errors and Model Uncertainty}

To assess practical robustness, we introduce three degradation
scenarios: (a) outdoor temperature forecast bias of $\pm\SI{3}{\kelvin}$;
(b) random forecast noise with $\sigma = \SI{2}{\kelvin}$; and (c)
$\pm\SI{25}{\percent}$ perturbation of the RC parameters $C_a$, $R_{as}$.
In all cases the disturbance observer compensates for the resulting
model mismatch, and comfort violations increase by at most three
intervals per week versus the nominal case. Energy consumption
increases by less than \SI{6}{\percent}, confirming that the closed-loop
system is robustly stable and practically reliable under realistic
operating conditions.

% =====================================================================
\section{Discussion}
\label{sec:discussion}
% =====================================================================

\textbf{Pre-conditioning and TOU exploitation.}
The most distinctive feature of MPC compared with reactive controllers
is its ability to exploit building thermal mass as a free energy buffer.
By pre-heating (or pre-cooling) the zone slightly in advance of the
occupied period, using off-peak electricity, MPC achieves cost
reductions of over 40\% relative to the thermostat. This load-shifting
capability is increasingly valuable as electricity grids integrate
higher shares of renewable generation, which introduce larger day-ahead
price volatility.

\textbf{Comfort--energy trade-off.}
The comfort weight $q$ in~\cref{eq:cost} determines the trade-off
between strict setpoint adherence and energy savings. Higher values
of $q$ reduce comfort violations at the expense of energy, while lower
values produce more aggressive energy savings but may allow the
temperature to approach constraint boundaries more
frequently~\cite{Serale2018}. In the current setting, $q = 100$
achieves four or fewer violations per seven days---corresponding to less
than \SI{1}{\hour} total discomfort---which satisfies typical occupant
acceptance criteria~\cite{ASHRAE55,Fanger1970}.

\textbf{Computational considerations.}
The QP~\eqref{eq:qp} has $N = 24$ decision variables and approximately
$5N = 120$ inequality constraints. On a commodity single-board computer
(Raspberry Pi 4, \SI{1.8}{\giga\hertz}), QUADPROG solves the problem
in under \SI{5}{\milli\second}, far below the \SI{15}{\minute} sampling
period. This confirms that MPC is computationally feasible for
deployment on embedded BMS hardware without dedicated
co-processors~\cite{Rawlings2017,Lofberg2004}.

\textbf{Model identification effort.}
The main practical barrier to deploying building MPC is the model
identification step. In our study, parameter identification required
approximately two weeks of logged operational data under normal HVAC
cycling. Bacher and Madsen~\cite{Bacher2011} showed that as few as
three to five days of data suffice for reliable RC parameter estimation
when the HVAC is operated in a persistently exciting
manner. Automated identification pipelines and physics-informed machine
learning models are active research directions that could reduce this
barrier further~\cite{Privara2013,Yang2015}.

\textbf{Limitations.}
This study addresses a single thermal zone. In multi-zone buildings,
inter-zone thermal coupling introduces coupling constraints that
complicate the MPC formulation~\cite{Morosan2010}. The RC model
also does not capture humidity or CO$_2$ concentration, which influence
occupant comfort and indoor air quality. Finally, the solar irradiance
and occupancy forecasts used here were assumed to be deterministic; a
stochastic or robust MPC formulation would improve worst-case guarantees
under uncertain forecasts~\cite{Oldewurtel2010}.

% =====================================================================
\section{Conclusion}
\label{sec:conclusion}
% =====================================================================

This paper presented a complete MPC framework for energy-efficient HVAC
control in a smart building automation context. A second-order RC
thermal network, identified from operational data, served as the
prediction model within a constrained quadratic optimization problem
solved over a six-hour receding horizon. A Luenberger disturbance
observer augmented the state with an estimated unmodeled heat flow,
compensating for the inevitable gap between the identified model and the
true building dynamics.

Simulation over seven-day winter and summer periods demonstrated that
MPC outperforms both a rule-based thermostat and a PI controller across
all evaluated metrics. Energy savings of \SI{31.2}{\percent} (winter)
and \SI{29.8}{\percent} (summer) were achieved versus the thermostat,
comfort violations were reduced by more than \SI{78}{\percent} versus
PI, and electricity cost was cut by over \SI{40}{\percent} through
time-of-use load shifting. The controller solved the underlying QP in
under \SI{5}{\milli\second} on embedded hardware, confirming practical
deployability.

Future work will extend the framework to multi-zone buildings with
thermal coupling, incorporate humidity and CO$_2$ into the comfort
formulation, and develop a stochastic MPC variant that formally handles
weather forecast uncertainty. Integration with on-site photovoltaic
generation and battery storage will also be investigated to further
reduce peak-grid demand and carbon footprint.

% =====================================================================
\section*{Acknowledgment}
% =====================================================================

The author thanks the anonymous reviewers for their helpful comments.
This work was supported in part by the National Natural Science
Foundation of China.

% =====================================================================
% Bibliography
% =====================================================================

\clearpage
\bibliographystyle{IEEEtranN}
\bibliography{paper2}

\ \\
\noindent All links were last verified on May~4, 2026.

\end{document}
